% Original flowchart for the owner's paper-reproduce workflow.
% Content: assets/decision-stages.mmd and the approved three choices (matrix cell, target, working mode).
% Renderer: LuaLaTeX/TikZ. Profile: schematic-designer Minimalist.
% Design lesson: aligned groups and reserved connector space, from the local
% lab-curriculum-flow tutorial. No example or upstream artwork/code copied.
% Palette: teal for the matrix cell; warm gold for the target; warm red for the mode. All alternatives
% within a choice share one color; labels retain the distinction without color.
% Geometry indicates choices and ordering only; there are no measured quantities.
% Generate minimalist-profile.sty from schematic-designer's shared profile.
\documentclass[tikz,border=2mm]{standalone}
\usepackage{minimalist-profile}
\usetikzlibrary{calc}
\colorlet{objectiveStroke}{msColor5!75!msInk}
\colorlet{objectiveText}{msColor6}
\colorlet{objectiveCard}{msColor5!8!msPaper}
\colorlet{objectiveGroup}{msColor5!3!msPaper}
\colorlet{objectiveBorder}{msColor5!35!msPaper}
\colorlet{depthStroke}{msColor3!60!msInk}
\colorlet{depthText}{msColor3!40!msInk}
\colorlet{depthCard}{msColor3!14!msPaper}
\colorlet{depthGroup}{msColor3!4!msPaper}
\colorlet{depthBorder}{msColor3!50!msPaper}
\colorlet{modeStroke}{msColor1!75!msInk}
\colorlet{modeText}{msColor1!70!msInk}
\colorlet{modeCard}{msColor1!7!msPaper}
\colorlet{modeGroup}{msColor1!3!msPaper}
\colorlet{modeBorder}{msColor1!30!msPaper}
\begin{document}
\pagecolor{msPaper}
\hyphenpenalty=10000
\exhyphenpenalty=10000
\begin{tikzpicture}[ms base,x=1mm,y=1mm,
    stage/.style={draw=msMuted!60,fill=msPaper,line width=\msLineContext,rounded corners=1.5mm},
    card/.style={draw=msInk,fill=msPaper,line width=\msLineStructure,rounded corners=1.4mm,minimum height=24mm,inner sep=0mm,outer sep=0mm},
    cell/.style={card,minimum width=62mm,minimum height=14mm},
    target/.style={card,minimum width=40mm},
    mode/.style={card,minimum width=52mm},
    heading/.style={ms title,anchor=west,inner sep=0mm},
    axis label/.style={font=\sffamily\fontsize{11pt}{13pt}\selectfont\bfseries,inner sep=0mm,text=objectiveText},
    card title/.style={ms title,anchor=north,inner sep=0mm,align=center},
    description/.style={ms body,anchor=north,inner sep=0mm,align=center}]

% Choice 1 is a two-by-two matrix: rows are the data, columns the analysis.
% The cells carry names only; the row and column labels define them.
\draw[stage,draw=objectiveBorder,fill=objectiveGroup] (0,0) rectangle (176,-58);
\node[heading,text=objectiveText] at (5,-6) {Choice 1: Which cell of the matrix?};
\node[axis label,anchor=south] at (68,-15.5) {Same analysis};
\node[axis label,anchor=south] at (134,-15.5) {Different analysis};
\node[axis label,anchor=east,align=right,text width=28mm] at (33,-27) {Same data};
\node[axis label,anchor=east,align=right,text width=28mm] at (33,-45) {Different data};
\node[cell,draw=objectiveStroke,fill=objectiveCard] (reproducible) at (68,-27) {};
\node[cell,draw=objectiveStroke,fill=objectiveCard] (robust) at (134,-27) {};
\node[cell,draw=objectiveStroke,fill=objectiveCard] (replicable) at (68,-45) {};
\node[cell,draw=objectiveStroke,fill=objectiveCard] (generalisable) at (134,-45) {};
\node[card title,anchor=center,text=objectiveText,text width=56mm] at (68,-27) {Reproducible};
\node[card title,anchor=center,text=objectiveText,text width=56mm] at (134,-27) {Robust};
\node[card title,anchor=center,text=objectiveText,text width=56mm] at (68,-45) {Replicable};
\node[card title,anchor=center,text=objectiveText,text width=56mm] at (134,-45) {Generalisable};

\begin{scope}[yshift=20mm]
\draw[ms arrow] (88,-78) -- (88,-92);

% Choice 2 is one group. No connector joins individual alternatives.
\draw[stage,draw=depthBorder,fill=depthGroup] (0,-92) rectangle (176,-134);
\node[heading,text=depthText] at (5,-98) {Choice 2: What is the target?};
\node[target,draw=depthStroke,fill=depthCard] (model) at (23,-117) {};
\node[target,draw=depthStroke,fill=depthCard] (claim) at (66.33,-117) {};
\node[target,draw=depthStroke,fill=depthCard] (result) at (109.67,-117) {};
\node[target,draw=depthStroke,fill=depthCard] (measurements) at (153,-117) {};
\node[card title,text=depthText,text width=36mm] at (23,-108) {Model or method};
\node[card title,text=depthText,text width=36mm] at (66.33,-108) {Claim or hypothesis};
\node[card title,text=depthText,text width=36mm] at (109.67,-108) {Result};
\node[card title,text=depthText,text width=36mm] at (153,-108) {Measurements};
\node[description,text width=36mm] at (23,-114) {Equations, parameters, conditions, and numerics, whole or in part.};
\node[description,text width=36mm] at (66.33,-114) {A claim, mechanism, or hypothesis posed as a neutral question.};
\node[description,text width=36mm] at (109.67,-114) {A figure, table, curve, or single value.};
\node[description,text width=36mm] at (153,-114) {Published experimental data, extracted or re-analysed.};

\draw[ms arrow] (88,-134) -- (88,-148);

% Choice 3 is the working mode: how much the researcher sees and decides.
\draw[stage,draw=modeBorder,fill=modeGroup] (0,-148) rectangle (176,-190);
\node[heading,text=modeText] at (5,-154) {Choice 3: How do you want to work?};
\node[mode,draw=modeStroke,fill=modeCard] (humanloop) at (30,-173) {};
\node[mode,draw=modeStroke,fill=modeCard] (autonomous) at (88,-173) {};
\node[mode,draw=modeStroke,fill=modeCard] (guided) at (146,-173) {};
\node[card title,text=modeText,text width=46mm] at (30,-164) {Human-in-the-loop (default)};
\node[card title,text=modeText,text width=46mm] at (88,-164) {Autonomous};
\node[card title,text=modeText,text width=46mm] at (146,-164) {Guided learning};
\node[description,text width=46mm] at (30,-170) {You make every scientific choice; the agent stops at each one.};
\node[description,text width=46mm] at (88,-170) {The agent works through the agreed plan and logs every choice it makes.};
\node[description,text width=46mm] at (146,-170) {The agent explains each step so you could redo it by hand.};

\node[ms small,anchor=north,inner sep=0mm,text=msMuted] at (88,-195) {Choose one cell, one target, and one mode; any combination is allowed. Definitions follow The Turing Way.};
\end{scope}
\end{tikzpicture}
\end{document}
